% from ega0-20-fr.tex
\oldpage[0\textsubscript{IV-1}]{143}
où l'opérateur bord $\partial$ est le composé
\begin{equation}
\tag{20.6.17.2}
\label{0.20.6.17.2-fr}
  \Upsilon_{C/B/\Lambda}\xrightarrow{s'}\Upsilon_{C/B/A}
  \longrightarrow\Omega^1_{B/A}\otimes_B C
\end{equation}
la seconde flèche étant l'injection canonique.
\end{theorem}

\begin{proof}
En écrivant la suite exacte d'homologie pour la suite exacte de complexes
\sref{0.20.6.16.1-fr}, on obtient \sref{0.20.6.17.1-fr}, l'homologie étant
nulle en degrés distincts de 0 et de 1 ; le fait que les homomorphismes de
cette suite exacte qui proviennent par fonctorialité de $j$ et de $p$ sont
bien ceux de l'énoncé est immédiat. Reste à vérifier que $\partial$ est égal
à \sref{0.20.6.17.2-fr} ; or un élément $z\in\Upsilon_{C/B/\Lambda}$ est
image par $p_1$ de $(0,z)$, d'où on déduit aussitôt que $\partial(z)$ est
l'image de $z$ par l'homomorphisme canonique
\[
  s_{B/A/\Lambda}\otimes1_C:
  \Omega^1_{B/\Lambda}\otimes_B C\to\Omega^1_{B/A}\otimes_B C.
\]
Notre assertion résulte de la commutativité du diagramme
\begin{equation}
\tag{20.6.17.3}
\label{0.20.6.17.3-fr}
\xymatrix{
  \Upsilon_{C/B/\Lambda} \ar[r]^{s'} \ar[dd] & \Upsilon_{C/B/A} \ar[dd] \\
  & \\
  \Omega^1_{B/\Lambda}\otimes_B C \ar[r] &
  \Omega^1_{B/A}\otimes_B C
}
\end{equation}
(cf.\ \sref{0.20.6.4.3-fr}).
\end{proof}

\begin{corollary}[20.6.18]
\label{0.20.6.18-fr}
\begin{enumerate}
  \item[\rm(i)] La suite de $C$-modules
\begin{equation}
\tag{20.6.18.1}
\label{0.20.6.18.1-fr}
  0\to\Upsilon^C_{B/A/\Lambda}\xrightarrow{v'}
  \Upsilon_{C/A/\Lambda}\xrightarrow{u'}
  \Upsilon_{C/B/\Lambda}\xrightarrow{s'}\Upsilon_{C/B/A}\to0
\end{equation}
est exacte.
  \item[\rm(ii)] Si $B$ est une $A$-algèbre formellement lisse relativement
  à $\Lambda$, on a $\Upsilon^C_{B/A/\Lambda}=0$.
\end{enumerate}
\end{corollary}

\begin{proof}
(i) Dans \sref{0.20.6.17.1-fr}, l'image de $\Upsilon_{C/B/\Lambda}$ par
$\partial$ est le noyau de $v_{C/B/A}$, donc $\Upsilon_{C/B/A}$ par
définition.

(ii) L'hypothèse entraîne que la suite
\[
  0\to\Omega^1_{A/\Lambda}\otimes_A B\to
  \Omega^1_{B/\Lambda}\to\Omega^1_{B/A}\to0
\]
est exacte et scindée \sref{0.20.5.7-fr} ; par tensorisation avec $C$, la
suite
\[
  0\to\Omega^1_{A/\Lambda}\otimes_A C\to
  \Omega^1_{B/\Lambda}\otimes_B C\to
  \Omega^1_{B/A}\otimes_B C\to0
\]
reste donc exacte, d'où notre assertion.
\end{proof}

\begin{corollary}[20.6.19]
\label{0.20.6.19-fr}
Soient $K$ un corps, $E$, $F$ deux extensions de $K$ telles que
$K\subset E\subset F$.
\begin{enumerate}
  \item[\rm(i)] Si $F$ est une extension séparable de $E$, on a
  $\Upsilon_{F/E/K}=0$ (autrement dit, l'homomorphisme canonique
  $\Omega^1_{E/K}\otimes_E F\to\Omega^1_{F/K}$ est injectif).
  \item[\rm(ii)] Inversement, si $F$ est une extension séparable de $K$ et
  si l'on a $\Upsilon_{F/E/K}=0$, alors $F$ est une extension séparable de
  $E$.
\end{enumerate}
\end{corollary}

\begin{proof}
Si $P$ est le corps premier de $K$, on a en effet la suite exacte
\sref{0.20.6.18-fr}
\[
  \Upsilon_{F/K/P}\to\Upsilon_{F/E/P}\to\Upsilon_{F/E/K}\to0.
\]
Si $\Upsilon_{F/E/P}=\Upsilon_{F/E}=0$, on a donc
$\Upsilon_{F/E/K}=0$, d'où (i) en vertu de \sref{0.20.6.3-fr} ;
inversement, si $\Upsilon_{F/E/K}=0$ et
$\Upsilon_{F/K/P}=\Upsilon_{F/K}=0$, on a
$\Upsilon_{F/E/P}=0$, d'où (ii) en vertu de \sref{0.20.6.3-fr}.
\end{proof}

